\documentclass[10pt,a4paper]{amsart}
\usepackage[margin=19mm]{geometry}
\usepackage{amsmath,amssymb,mathtools}
\usepackage[scr=rsfs,cal=euler]{mathalfa}
\usepackage{xfrac}
\usepackage[unicode,hidelinks,bookmarks=false]{hyperref}
\newcommand{\ie}{i.e.\ }
\newcommand{\cf}{cf.\ }
\newcommand{\resp}{respectively\ }
\newcommand{\Subsection}{\subsection}
\providecommand{\nobreakdash}{\nobreak\hbox{-}}
\setlength{\emergencystretch}{2em}
\multlinegap0pt
\usepackage{mathtools}
\usepackage{amssymb}
\usepackage{bm}
\usepackage{enumerate}
\usepackage[normalem]{ulem}
\usepackage{tikz}
\usepackage{xcolor}
\usepackage{cancel}
\newtheorem{lemma}{Lemma}
\def\A{{\mathcal A}}
\def\R{{\mathbb R}}
\title{Quantitative propagation of chaos for non-exchangeable diffusions via first-passage percolation}
\author{Daniel Lacker, Lane Chun Yeung, Fuzhong Zhou}
\date{Source: arXiv:2409.08882v2 (CC BY 4.0)}
\begin{document}
\maketitle

\noindent\emph{Source and licence.} Quantitative propagation of chaos for non-exchangeable diffusions via first-passage percolation. Version arXiv:2409.08882v2 (CC BY 4.0), distributed by arXiv under the Creative Commons Attribution 4.0 International licence, http://creativecommons.org/licenses/by/4.0/, as recorded in the arXiv licence metadata for this exact version and shown on the abstract page https://arxiv.org/abs/2409.08882v2. Full licence notice: \texttt{licenses/arxiv-2409.08882-CC-BY-4.0.txt}.\par\smallskip

\noindent\textit{Editorial scope.} Counted unit: the article's lemma le:AF-quadratic (source0.tex lines 1819--1825). The article's complete original proof of it is delivered with it (source0.tex lines 1826--1860, one \texttt{proof} environment of 35 lines). No mathematics is written, repaired or re-derived: the statement and the proof are the article's own lines, in the article's own notation, and no retained line is substituted or re-flowed.  Retained uncounted as the source-local context the proof needs: lines 1636--1638; lines 1671--1673; lines 1681--1686; lines 1733--1739.  Nothing was omitted from any retained span, so the package carries no omission record.  No retained line contains a citation, so the wrapper carries no bibliography and no bibliography entry.  No retained line contains a figure environment, an image-inclusion command or drawing code, so no image is delivered and the package declares no asset dependency.  Editorial formatting only, in addition: an AMS \texttt{amsart} preamble that restates the article's own declarations and macros used by the retained lines, namely the environments \texttt{lemma} on separate counters, together with the standard packages those lines need.  The retained environments print in this document as Lemma 1 in order of appearance, whereas the article prints the same items with the numbering of its own full document.  The complete native source of the article as distributed in the arXiv e-print for this exact version is bundled unchanged as \texttt{sources/165215/source0.tex}.
			\begin{equation} \tag{R-rows}
		\max_{1\leq i\leq n} \sum_{j=1}^n R_{ij} \leq 1. \label{cond:R row.sum}
	\end{equation}  Then we have the following estimates, for any $t \ge 0$ and $v \subset [n]$:

\begin{equation}
\A^{R} F(v) = \frac{2\gamma}{\sigma^2}\sum_{j \notin v}\Big(\sum_{i \in v} R_{ij}\Big) (F(vj)-F(v)), \label{eq:AF-basic}
\end{equation}

\begin{lemma} \label{le:AF-polynomial}
Let $\ell \ge 1$. Then, for $v \subset [n]$,
\begin{align}
		\A^{R}[|v|^\ell] \le \frac{2\gamma}{\sigma^2}   |v|\big((|v|+1)^{\ell}-|v|^\ell\big). \label{pf:polynomial1}
		\end{align}
\end{lemma}

\begin{lemma} \label{le:AF-linear}
Let $x \in \R^n$ have nonnegative entries, and let $\ell \ge 0$ be an integer. Let $v \subset [n]$. 
\begin{equation}
\A^{R}[|v|^\ell \langle 1_v,x\rangle ] 
 \le \frac{2\gamma}{\sigma^2}(|v|+1)^\ell \langle 1_v, R x\rangle + \frac{2\gamma}{\sigma^2} |v| \big((|v|+1)^{\ell} - |v|^\ell\big)\langle 1_v, x\rangle . \label{pf:linear1}
\end{equation}
\end{lemma}

\begin{lemma} \label{le:AF-quadratic}
	Let $ G $ be  an $n \times n$ matrix with nonnegative entries. Let $v \subset [n]$. Then
		\begin{align} 
			\A^{R} [\langle 1_v, G1_v\rangle ] &\le \frac{2\gamma}{\sigma^2} \langle 1_v , R G_{\mathrm{diag}}\rangle  + \frac{2\gamma}{\sigma^2}   \langle 1_v, \big( R G +GR^\top \big) 1_v\rangle , \label{eq: AF.bound.0} \\
			\A^{R} [  |v| \langle 1_v, G1_v\rangle ]  &\le \frac{2\gamma}{\sigma^2} (|v|+ 1) \Big[  \langle 1_v , R G_{\mathrm{diag}}\rangle  + \langle 1_v, \big( R G +GR^\top  \big) 1_v\rangle\Big] + \frac{2\gamma}{\sigma^2}  |v|\langle 1_v,G1_v\rangle .  \label{eq: AF.bound.1}
		\end{align} 
\end{lemma}

\begin{proof} 
As in the proofs of Lemmas \ref{le:AF-polynomial} and \ref{le:AF-linear}, we assume without loss of generality that $\sigma = \sqrt{2}$. We start with \eqref{eq: AF.bound.0}. Let $F(v) = \langle 1_v, G1_v\rangle = \sum_{i,r \in v}G_{ir}$. 	For $j \notin v$ we compute
		\begin{align*}
			F(vj) - F(v) &= \sum_{i,r \in vj}G_{ir} - \sum_{i,r \in v}G_{ir} = G_{jj} + \sum_{r \in v}(G_{rj} + G_{jr}).
		\end{align*}
		Thus, using \eqref{eq:AF-basic} and the nonnegativity of the entries of $R$,
		\begin{align*}
			\A^{R} F(v) &=  \gamma \sum_{j \notin v}\bigg(\sum_{i \in v} R_{ij}\bigg)\bigg(G_{jj} + \sum_{r \in v}(G_{rj} + G_{jr})\bigg) \\
			&\le \gamma \sum_{i \in v}\sum_{j=1}^n R_{ij}G_{jj} + \gamma \sum_{i,r \in v} \sum_{j =1}^n (R_{ij}G_{rj} + R_{ij}G_{jr}) \\
			&= \gamma \langle 1_v , R G_{\mathrm{diag}}\rangle  + \gamma  \langle 1_v, \big( R G +GR^\top \big) 1_v\rangle .
		\end{align*} 
	We next turn to \eqref{eq: AF.bound.1}. Set $F(v)=|v|\langle 1_v, G1_v\rangle$. For $j \notin v$ we compute  
		\begin{align*}
			F(vj) - F(v) &= (|v| +1)\sum_{\ell, r\in vj}G_{\ell r} -  |v| \sum_{\ell, r\in v}G_{\ell r} \\
			&= |v| \bigg(G_{jj} + \sum_{r \in v}(G_{rj} + G_{jr})\bigg) + \sum_{\ell, r\in vj} G_{\ell r}.
		\end{align*}
		Thus
		\begin{align*}
			\A^{R}  F(v) &= \gamma \sum_{j \notin v}\bigg(\sum_{i \in v} R_{ij}\bigg)\bigg(  |v| \bigg(G_{jj} + \sum_{r \in v}(G_{rj} + G_{jr})\bigg) + \sum_{\ell, r\in vj} G_{\ell r} \bigg) \\
			&\le \gamma  |v|  \sum_{i \in v}\sum_{j=1}^n R_{ij}G_{jj} + \gamma  |v| \sum_{i,r \in v} \sum_{j =1}^n \Big(R_{ij}G_{rj} + R_{ij}G_{jr}\Big) 
			+ \gamma \sum_{i\in v}\sum_{j \notin  v} R_{ij} \sum_{\ell, r\in vj} G_{\ell r}  \\
			&= \gamma |v|  \langle 1_v, R G_{\mathrm{diag}}\rangle  + \gamma |v|  \langle 1_v, \big( R G +GR^\top \big) 1_v\rangle + \gamma \sum_{i\in v}\sum_{j\notin  v} R_{ij} \sum_{\ell, r\in vj} G_{\ell r}.
		\end{align*}
		We simplify the last term by splitting the sum into four cases, depending on whether $\ell$ and $r$ equal $j$, or both, or neither:
		\begin{align*}
			\sum_{i\in v}\sum_{j\notin v} R_{ij} \sum_{\ell, r\in vj} G_{\ell r} 
			&\le \sum_{i\in v} \bigg( \sum_{\ell,r\in v} G_{\ell r} + \sum_{\ell\in v} \sum_{j=1}^n R_{ij}G_{\ell j} + \sum_{r\in v} \sum_{j=1}^nR_{ij}G_{j r} + \sum_{j=1}^nR_{ij}G_{j j} \bigg) \\
			&=  |v| \langle 1_v, G 1_v \rangle + \langle 1_v, GR^\top 1_v\rangle  + \langle 1_v, R G 1_v \rangle + \langle 1_v,R G_{\mathrm{diag}}\rangle,
		\end{align*}
		where we used the assumption  \eqref{cond:R row.sum} to remove the $R$ in the first step, and we used nonnegativity of the entries of $R$ and $G$ throughout. 
		Combining this with the previous inequality yields
		\begin{align*}
			\A^{R} F(v) &\le \gamma(|v|+ 1) \langle 1_v,R G_{\mathrm{diag}}\rangle + \gamma (|v|+ 1) \langle 1_v , \big( R G +GR^\top \big) 1_v\rangle + \gamma|v| \langle 1_v, G 1_v \rangle. \qedhere
		\end{align*} 	 
\end{proof}

\end{document}
